\section{Client--Bridge Wire Format}
\label{app:wire-format}
% ================================================================================

This appendix is \emph{normative} for the reference transport: a bridge
that claims to speak it and emits or accepts messages in a different
shape is non-conformant at the wire level, even if its committed event
log and final session state agree with the contract. A binding to
another transport keeps only the parts named in
Section~\ref{sec:wire-binding}. The conformance unit remains the committed event log
and the final session state (Section~\ref{sec:formal-status}); the
wire format is what makes byte-level interoperability between
bridges achievable.

\paragraph{What is exercised.} The normative status is wider than the
evidence, and the evidence is stated here. The signed byte strings, the
reason names and the canonical encoding are pinned by known-answer vectors
and checked by both bridges. The envelope is exercised end to end on both
reference implementations: a wire front-end in each
(\texttt{supplementary/wire\_server.py} and
\texttt{supplementary/rust-bridge/src/wire\_server.rs}) decodes a request,
refuses bytes that are not canonical JSON, dispatches the op-types
\texttt{attach}, \texttt{challenge}, \texttt{add\_client},
\texttt{reattach}, \texttt{join}, \texttt{disconnect},
\texttt{escalate}, \texttt{resolve}, \texttt{submit\_tool\_call},
\texttt{answer\_tool\_call}, \texttt{delivered}, \texttt{completed},
\texttt{narrow}, \texttt{extend} and \texttt{terminate}, and the operator
events \texttt{revoke}, \texttt{add\_key} and \texttt{operator\_terminate},
to its bridge, and returns the response. The vectors run through each with
real signatures that the bridge verifies from the message, and the session
of each vector is created by an \texttt{attach} message
(\texttt{run\_vectors.py -{}-wire} and the Rust test
\texttt{all\_vectors\_pass\_through\_the\_wire\_server}). The two
front-ends share the reason names and the hook for the transport's
caller check, and were written by the same author. What the vectors do not
send as messages are the clock and restart hooks and the tampering of logs
and registries, which are test-harness actions, and a re-attach left
incomplete or flagged as a transport replay.

The wire format is JSON, encoded canonically per RFC~8785~\cite{rundgren2020json}
for the value types the contract uses: strings (any Unicode text, emitted
as raw UTF-8 with the RFC~8785 escapes), integers whose magnitude is at
most $2^{53}$, finite numbers (ECMAScript formatting, so a timestamp
\texttt{1.0} is encoded \texttt{1}), arrays, and objects with string
member names sorted by UTF-16 code units. A value outside that set (NaN,
an infinity, an integer beyond $2^{53}$) has no canonical form and a
bridge MUST refuse to hash or sign it. Both reference implementations
produce the same bytes, pinned by known-answer vectors
(Section~\ref{sec:wire-kat}).
A future revision MAY define a binary encoding (e.g., protobuf or
CBOR) as an alternative binding, provided it preserves a canonical
byte representation for signatures.

\subsection{What a transport binding must keep}
\label{sec:wire-binding}

A deployment may carry the contract in its own messages, for example
inside the frames of an existing agent protocol, instead of in the
envelope below. Such a binding is conformant at the wire level when it
keeps the parts listed here and may change the rest.
\begin{itemize}[leftmargin=2em]
\item \textbf{Signature scope.} The arrays in the ``Signature over''
column and in the signature-scope subsection, with their
domain-separation tags and the order of their elements, are normative.
A client and a bridge that sign and verify other bytes do not
interoperate with a bridge that follows the vectors.
\item \textbf{Reason names.} The reasons in Table~\ref{tab:wire-reasons}
are the rejection vocabulary of the contract. A binding MAY wrap a
reason in its own error structure, but it MUST preserve the reason name
so that a client can tell \texttt{prefix\_mismatch} from
\texttt{nonce\_used}.
\item \textbf{Canonical encoding of signed bytes.} What is signed is the
RFC~8785 canonical JSON of the listed array, whatever encoding carries
the message. A binding that does not use JSON on the wire still
reconstructs these bytes before it signs or verifies.
\end{itemize}
The envelope with its \texttt{op}, \texttt{status} and \texttt{reason}
fields, the one-request-one-response framing and the op-type names are
the reference transport only. A binding MAY map each op-type onto its
own message type, and the vectors then act as the acceptance interface
at the level of events and final state, not of message shape. The
requirements on unsigned ops (authentication and freshness by the
transport) apply to every binding.

\subsection{Envelope}

Every client$\to$bridge request is a JSON object with an \texttt{op}
field naming the op-type and a \texttt{session\_id} field (except
\texttt{restart}, which is bridge-level). Every bridge$\to$client
response is a JSON object with a \texttt{status} field and a
\texttt{session\_id} field. Errors use
\texttt{status: "rejected"} and a short snake\_case \texttt{reason}.

\begin{lstlisting}
// client -> bridge
{"op": "<op-type>", "session_id": "<sid>", ...op-specific fields...}

// bridge -> client (success)
{"status": "<status>", "session_id": "<sid>", ...response fields...}

// bridge -> client (error)
{"status": "rejected", "session_id": "<sid>", "reason": "<reason>"}
\end{lstlisting}

\subsection{Op-types}

Table~\ref{tab:wire-ops} lists the op-types. The
\texttt{signature} field, where present, is a hex-encoded signature (Ed25519 or ECDSA P-256 in the
reference implementations), over the canonical JSON of the array listed in the
``Signature over'' column. The bridge verifies signatures against
the registry of client public keys (I11).

\footnotesize
\begin{longtable}{p{0.13\columnwidth} p{0.26\columnwidth} p{0.26\columnwidth} p{0.21\columnwidth}}
\caption{Client--bridge op-types.}\label{tab:wire-ops}\\
\toprule
\textbf{Op} & \textbf{Request fields} & \textbf{Response} & \textbf{Signature over} \\
\midrule
\endfirsthead
\toprule
\textbf{Op} & \textbf{Request fields} & \textbf{Response} & \textbf{Signature over} \\
\midrule
\endhead
\texttt{attach} &
\texttt{session\_id}, \texttt{principal}, \texttt{scope[]}, \texttt{lineage[]}, \texttt{key\_id}, \texttt{scheme}, \texttt{public\_key}, \texttt{grant}, \texttt{signature} &
\texttt{status: "attached"}, \texttt{seq: 0}, \texttt{hash} \emph{or} \texttt{status: "rejected"}, \texttt{reason} &
\texttt{["bc/v1/attach", bridge\_id, sid, principal, scope, lineage, key\_id, scheme, public\_key]} \\
\texttt{add\_client} &
\texttt{session\_id}, \texttt{via\_key\_id}, \texttt{new\_key\_id}, \texttt{scheme}, \texttt{new\_public\_key}, \texttt{nonce}, \texttt{signature} &
\texttt{status: "added"} \emph{or} \texttt{status: "rejected"}, \texttt{reason} &
\texttt{["bc/v1/add-client", sid, via\_key\_id, new\_key\_id, scheme, new\_public\_key, nonce]} \\
\texttt{challenge} &
\texttt{session\_id}, \texttt{key\_id} &
\texttt{status: "nonce"}, \texttt{nonce}, \texttt{ttl\_seconds} &
--- (no signature) \\
\texttt{reattach} &
\texttt{session\_id}, \texttt{key\_id}, \texttt{nonce}, \texttt{seq}, \texttt{hash}, \texttt{signature} &
\texttt{status: "verified"}, \texttt{events[]} \emph{or} \texttt{status: "rejected"}, \texttt{reason} &
\texttt{["bc/v1/reattach", sid, key\_id, nonce, seq, hash]} \\
\texttt{join} &
same as \texttt{reattach} (session is attached) &
\texttt{status: "joined"}, \texttt{events[]} &
\texttt{["bc/v1/join", sid, key\_id, nonce, seq, hash]} \\
\texttt{disconnect} &
\texttt{session\_id}, \texttt{key\_id}? (absent = all), \texttt{reason} &
\texttt{status: "disconnected"} &
--- \\
\texttt{escalate} &
\texttt{session\_id}, \texttt{trigger}? &
\texttt{status: "escalated"} &
--- \\
\texttt{resolve} &
\texttt{session\_id} &
\texttt{status: "resolved"} &
--- \\
\texttt{submit\_tool\_call} &
\texttt{session\_id}, \texttt{tool}, \texttt{key}?, \texttt{permission}?, \texttt{content} &
\texttt{status: "queued"} or \texttt{"forwarded"} (under the forward profile), with the call's \texttt{key}; a duplicate key or a full queue is refused through the error envelope (\texttt{duplicate\_key}, \texttt{queue\_full}) &
--- \\
\texttt{answer\_tool\_call} &
\texttt{session\_id}, \texttt{key}, \texttt{shown\_hash}, \texttt{decision} &
\texttt{status: "answered"} &
--- \\
\texttt{delivered} &
\texttt{session\_id}, \texttt{key} &
\texttt{status: "delivered"} &
--- \\
\texttt{completed} &
\texttt{session\_id}, \texttt{key}, \texttt{result\_hash}? &
\texttt{status: "completed"} &
--- \\
\texttt{narrow} &
\texttt{session\_id}, \texttt{new\_scope[]} &
\texttt{status: "narrowed"|"rejected"}, \texttt{new\_version}?, \texttt{reason}? &
--- \\
\texttt{extend} &
\texttt{session\_id}, \texttt{new\_scope[]}, \texttt{source} &
\texttt{status: "rejected"}, \texttt{reason: "extension\_rejected"},
\texttt{current\_scope[]} for a widening asked for (I3); \texttt{status:
"revert\_requested"} for a widening the runtime announced
(\texttt{source: "observed"}), recorded and to be reverted;
\texttt{status: "unchanged"} for a subset; all three carry
\texttt{current\_scope[]}, the scope in force &
--- \\
\texttt{terminate} &
\texttt{session\_id}, \texttt{reason} &
\texttt{status: "terminated"} &
--- (refused as \texttt{key\_exhausted} on a session without active keys) \\
\texttt{revoke} &
\texttt{session\_id}, \texttt{key\_id}, \texttt{reason}, \texttt{signature}?, \texttt{requested\_by}? &
\texttt{status: "revoked"|"requested"} &
\texttt{["bc/v1/revoke", bridge\_id, sid, key\_id, reason, seq, head]} (operator) \\
\texttt{add\_key} &
\texttt{session\_id}, \texttt{key\_id}, \texttt{scheme}, \texttt{public\_key}, \texttt{signature} &
\texttt{status: "key\_added"} &
\texttt{["bc/v1/add-key", bridge\_id, sid, key\_id, scheme, public\_key, seq, head]} (operator) \\
\texttt{operator\_terminate} &
\texttt{session\_id}, \texttt{reason}, \texttt{signature} &
\texttt{status: "terminated"} &
\texttt{["bc/v1/terminate", bridge\_id, sid, reason, seq, head]} (operator) \\
\bottomrule
\end{longtable}

\noindent The op-types \texttt{tick} (timer evaluation) and
\texttt{restart} (bridge recovery) are test-harness and
operator actions respectively, not client--bridge wire ops. They
are carried by the operator channel, not by the client--bridge wire
format defined here.

\subsection{Signature scope}

In every array, \texttt{sid} is the session identifier. Where
\texttt{bridge\_id} appears it is the identity of the bridge that verifies the
request; it is deployment configuration known to both sides and not a field of
the request. The arrays of \texttt{add\_client}, \texttt{reattach} and
\texttt{join} do not repeat it: each carries a nonce that this bridge issued
for this session and this key, and the nonce is spent by the first authentic
presentation, so it binds the request to the bridge that issued it.

These op-types carry a \texttt{signature} field:
\begin{itemize}[leftmargin=2em]
\item \texttt{attach}: the client proves it holds the private key of
the \texttt{public\_key} it is registering. The signature is made with
that key over the canonical JSON of the array
\texttt{["bc/v1/attach", bridge\_id, sid, principal, scope, lineage,
key\_id, scheme, public\_key]}, where \texttt{bridge\_id} is the identity
of the bridge the request is addressed to (deployment configuration known
to both sides; it is not a field of the request), \texttt{scheme} is the
scheme tag and \texttt{public\_key} the hex string registered. The array
has the fields of the grant (below) under a different tag, so the proof
and the grant cannot stand for one another. A missing proof, or one that
does not verify under the presented key, is answered with the reason
\texttt{possession\_invalid}; nothing is committed and no key is
registered.
\item \texttt{add\_client}: the connected client (identified by
\texttt{via\_key\_id}) authorises the new key. The signature is over
\texttt{["bc/v1/add-client", sid, via\_key\_id, new\_key\_id, scheme,
new\_public\_key, nonce]}, where \texttt{nonce} was issued by a
\texttt{challenge} naming \texttt{via\_key\_id} and is spent by the first
authentic presentation, so the authorisation cannot be replayed, and
\texttt{new\_public\_key} is the hex string registered. The bridge verifies
under the registered key of \texttt{via\_key\_id}, which must be connected
and not revoked. The refusals are \texttt{via\_not\_connected},
\texttt{bad\_signature} (missing, forged, or made for another operation,
session or key), \texttt{revoked\_key}, \texttt{nonce\_unknown},
\texttt{nonce\_used}, \texttt{nonce\_expired} and
\texttt{key\_already\_registered}. The signature is verified first: a caller
whose signature does not verify is answered \texttt{bad\_signature} whatever
the state of the session, and only a caller whose signature verified is told
that the session is \texttt{unknown\_session}, \texttt{terminated},
\texttt{quarantined} or \texttt{key\_exhausted}. A refusal that comes
before the signature has verified commits no event, so a caller that holds no
registered key writes nothing to the log.
\item \texttt{reattach} and \texttt{join}: the client proves
possession of a registered key and binds the request to the
challenge nonce and the last event it saw. The signature is over
\texttt{["bc/v1/reattach", sid, key\_id, nonce, seq, hash]}
(for \texttt{reattach}) or
\texttt{["bc/v1/join", sid, key\_id, nonce, seq, hash]}
(for \texttt{join}); this is the one definition, and the specification
and the state machine refer to it. The domain-separation tag
(\texttt{"bc/v1/reattach"} vs \texttt{"bc/v1/join"})
ensures that a \texttt{reattach} and a \texttt{join} signing the
same underlying values produce different signatures, preventing
cross-operation replay, and the key id ties the proof to the identity
that signs it. The bridge verifies against the operation it received,
not one the request names, so a proof made for the other operation is
refused as \texttt{bad\_signature}. A \texttt{reattach} for a session
that is attached is refused as \texttt{already\_attached}, and a
\texttt{join} for one that is not attached is refused.
\end{itemize}
The reference signature scheme is Ed25519, hex-encoded. A deployment
MAY use another scheme that is existentially unforgeable (A3), for
example ECDSA P-256 where a client platform lacks Ed25519; the
deployment's verifier then fixes the encoding of signature and public
key, and the signed bytes are unchanged. The reference implementations
also implement ECDSA P-256 with SHA-256 (public key: the 65-byte SEC1
uncompressed point; signature: the 64-byte $r \Vert s$ form, as
WebCrypto produces it), and the registry records a scheme tag with each
key. Vectors V59--V62 and V94 use P-256 (the file \texttt{p256\_kat.json} fixes its bytes); the other vectors use Ed25519. The
bridge verifies against
the registry of client public keys (I11); the registry holds
verification material only.

\paragraph{Known-answer vectors.}\label{sec:wire-kat}
\texttt{vectors/bytes\_kat.json} fixes the bytes, not only the behaviour.
It lists, with the bytes as text and as hex: the canonical form of
inputs that cover non-ASCII strings, escapes, key order beyond the
basic plane, integral and fractional numbers and exponent forms; inputs
that have no canonical form; the canonical bytes and hash of an event
and of tool-call content (so the event hash and the content hash are
pinned, including a non-ASCII payload and an integral timestamp); and
re-attach and join proofs under a fixed Ed25519 seed, with the exact
signed bytes and the signature (Ed25519 is deterministic, so the
signature is fixed), together with signatures that must not verify (the
other operation, another key id, the untagged tuple, the tuple without
the key id). Both reference implementations read the file, and the
canonical forms and signatures were also reproduced outside both, with
the string and number rules of JavaScript \texttt{JSON.stringify} and the
Ed25519 verifier of Node.js (OpenSSL). An independent implementation confirms
byte-level conformance by reproducing the file.

\paragraph{Authentication of unsigned ops.}
The op-types \texttt{resolve}, \texttt{answer\_tool\_call},
\texttt{escalate}, \texttt{delivered}, \texttt{completed},
\texttt{narrow}, \texttt{extend}, \texttt{disconnect},
\texttt{submit\_tool\_call}, and \texttt{terminate} do not
carry a client signature in the wire format. They MUST be
authenticated by the transport binding (e.g., TLS client
certificates, mTLS, or a session token issued at \texttt{attach}).
The bridge MUST verify the caller identity and reject unsigned
ops from unauthenticated transports. Unauthenticated answers are a
deployment risk: anyone who can speak the wire can inject a
\texttt{resolve} or \texttt{answer\_tool\_call}. The deployment
MUST document which transport-level authentication is used and
under which identity these ops are accepted. The transport binding
also carries freshness for these ops: they have no nonce or
\texttt{seq} of their own, so the contract does not make them
replay-safe on a transport that does not. Signing them is not
required and is not specified: an approval is bound to the request it
answers by a hash (I14), not to a signing identity, so who approved is
exactly what the transport authenticated, no more. The unsigned
\texttt{terminate} is the exception that the bridge itself narrows: on a
session without active keys it is refused (\texttt{key\_exhausted}) and only
the operator's signed termination ends the session.

\paragraph{Grant.}
The \texttt{grant} field of \texttt{attach} carries the scope grant that
the bridge presents to the deployment's grant verifier before it commits
\texttt{ATTACHED} (Section~\ref{sec:spec}, authorisation boundary). It is
an opaque string to the contract and is not covered by the client
signature: it authenticates itself, because it binds the bridge identity,
the session identifier, the principal, the scope, the lineage and the
client key (identifier, scheme tag and public key) under the domain tag
\texttt{"bc/v1/grant"}. Binding the public key stops a party that sees
a grant in transit from attaching the session under the granted key
identifier with a public key of its own; binding the bridge identity
stops a grant from being replayed at another bridge. The reference
verifier takes a hex Ed25519 signature by the operator over the
canonical JSON of \texttt{["bc/v1/grant", bridge\_id, sid, principal,
scope, lineage, key\_id, scheme, public\_key]}; a
deployment may use another format, such as an AGTP identity document or
a mission contract. A missing grant, a grant that does not verify, or a
bridge with no verifier is answered with the reason
\texttt{grant\_invalid}.

\paragraph{Operator signatures.}
A revocation, a key addition and a termination each carry a hex Ed25519
signature by the operator key, the key that signs scope grants. The signed
bytes are the canonical JSON of
\texttt{["bc/v1/revoke", bridge\_id, sid, key\_id, reason, seq, head]},
\texttt{["bc/v1/add-key", bridge\_id, sid, key\_id, scheme, public\_key,
seq, head]} and
\texttt{["bc/v1/terminate", bridge\_id, sid, reason, seq, head]}, where
\texttt{bridge\_id} is the identity of the bridge that verifies (the one
the grant binds), \texttt{public\_key} is the hex string registered,
\texttt{scheme} its scheme tag, and \texttt{seq} and \texttt{head} are the
sequence number and the event hash of the last committed event of the
session at the moment of signing. The bridge verifies against its current
last event, so a signature holds only while nothing has been committed
after that position and only at the bridge it names. The domain
tags keep these signatures, and the grant, from standing for one another.
A missing signature, one that does not verify, one by another key, or one
over another operation, session, bridge or log position is answered with
the reason \texttt{operator\_signature\_invalid}, and nothing is committed. A
revocation request without a signature (\texttt{KEY\_REVOCATION\_REQUESTED})
is recorded as a request and revokes nothing.

\subsection{Error reasons}

All rejections use \texttt{status: "rejected"} with a short
snake\_case \texttt{reason}. The defined reasons are:

\footnotesize
\begin{longtable}{p{0.28\columnwidth} p{0.60\columnwidth}}
\caption{Error reasons.}\label{tab:wire-reasons}\\
\toprule
\textbf{Reason} & \textbf{Meaning} \\
\midrule
\endfirsthead
\toprule
\textbf{Reason} & \textbf{Meaning} \\
\midrule
\endhead
\texttt{unknown\_session} & The \texttt{session\_id} does not name a known session. \\
\texttt{unknown\_key} & The \texttt{key\_id} is not in the session's registry $K(s)$. A re-attach or join naming such a key is recorded with this reason but answered \texttt{bad\_signature}, so a caller learns nothing about which key ids a session registers (I11); \texttt{revoke} names it. \\
\texttt{nonce\_used} & The nonce has already been spent (I11d). \\
\texttt{nonce\_expired} & The nonce is older than $T_{\text{nonce}}$ (I9). \\
\texttt{nonce\_unknown} & The nonce was never issued by the bridge for this session (I11b). \\
\texttt{nonce\_binding} & The nonce was issued for another key id, or for this one before it was registered or while it was revoked (I11b). Recorded only; the request is answered \texttt{bad\_signature} and the nonce is not spent. \\
\texttt{prefix\_mismatch} & The $(\mathit{seq}, \mathit{hash})$ pair is not a prefix of the bridge's chain (I4). \\
\texttt{resume\_in\_progress} & A re-attach is already being verified (I11d). \\
\texttt{already\_attached} & An authentic \texttt{reattach} reached a session that is attached; \texttt{join} is the only way into an attached session. Given only after the signature, the revocation, the binding and the nonce have been checked; the nonce is spent. \\
\texttt{not\_attached} & An authentic \texttt{join} reached a session that is not attached. Given only after the same checks; the nonce is spent. \\
\texttt{transport\_replay} & The transport reported a 0-RTT replay (I8). \\
\texttt{queue\_full} & The tool-call queue is at $Q_{\max}$ (I13). \\
\texttt{unanswered} & The call has no answer at drain (I14). \\
\texttt{denied} & The call was answered with \texttt{deny} (I14). \\
\texttt{answer\_mismatch} & $h_a \neq h_q$ at drain (I14). \\
\texttt{scope} & The call's permission is not in the current scope, at drain (I13) or, for a call submitted with a client attached and no escalation pending, at submission (I3). \\
\texttt{call\_ttl} & The queued call outlived $T_{\text{call}}$ (I13). \\
\texttt{approval\_ttl} & The approval outlived $T_{\text{approve}}$ (I13). \\
\texttt{terminated} & The session is in $\bot$ (I6). \\
\texttt{quarantined} & The session is quarantined (rule~Q). \\
\texttt{not\_a\_strict\_narrowing} & The proposed scope equals or widens the current scope (I10). \\
\texttt{refused} & The uniform refusal that names no cause: an unknown, terminated or quarantined session at re-attach or join (I11a), and the answer to a request that finds its session ended by a timer that expired in an earlier tick. \\
\texttt{unsupported\_op} & The reference wire server does not serve the op-type named. \\
\texttt{bad\_signature} & The signature does not verify, is malformed, or (re-attach and join) names a key id the session does not register. \\
\texttt{via\_not\_connected} & \texttt{add\_client} names an authorising key that is not connected (I15). \\
\texttt{key\_already\_registered} & \texttt{add\_client} names a key id already in $K(s)$. \\
\texttt{bad\_canonical} & The signed bytes are not canonical JSON per RFC 8785. \\
\texttt{revoked\_key} & The key has been revoked (I11-REVOKE); re-attach is rejected. \\
\texttt{key\_exhausted} & The session has no active keys (I11-REVOKE); all client operations are rejected, as is an unsigned \texttt{terminate}. \\
\texttt{not\_key\_exhausted} & \texttt{add\_key} or \texttt{operator\_terminate} reached a session that has active keys. \\
\texttt{extension\_rejected} & \texttt{extend} asked for a scope wider than the current one (I3); the response carries \texttt{current\_scope}. \\
\texttt{duplicate\_key} & A tool call with this idempotency key is already queued (I13). \\
\texttt{grant\_invalid} & The scope grant is missing or does not verify for this bridge, session, scope and client key (authorisation boundary). \\
\texttt{unauthenticated} & The transport binding did not authenticate the caller of an unsigned op (the paragraph on unsigned ops); nothing is committed. \\
\texttt{possession\_invalid} & The attach carries no proof of possession of the registered key, or the proof does not verify under it; nothing is committed and no key is registered. \\
\texttt{operator\_signature\_invalid} & The operator signature on a revocation, key addition or termination is missing, does not verify under the operator key, or covers another operation, bridge or log position (I-OPERATOR); nothing is committed. \\
\bottomrule
\end{longtable}

\subsection{Worked example}

A challenge--reattach sequence:

\begin{lstlisting}
-> {"op": "challenge", "session_id": "sess-001",
    "key_id": "ck-7f3a"}
<- {"status": "nonce", "session_id": "sess-001",
    "nonce": "9c1f5e0a77d24b3e8a6f21d0c4b7e913",
    "ttl_seconds": 60}

-> {"op": "reattach", "session_id": "sess-001",
    "key_id": "ck-7f3a",
    "nonce": "9c1f5e0a77d24b3e8a6f21d0c4b7e913",
    "seq": 6,
    "hash": "a7ffc6f8bf1ed76651c14756a061d662f3f7...",
    "signature": "q3Zk0f..."}
<- {"status": "verified", "session_id": "sess-001",
    "events": [{"sequence": 7, ...}, {"sequence": 8, ...}]}
\end{lstlisting}

\noindent A rejection:

\begin{lstlisting}
-> {"op": "reattach", "session_id": "sess-001",
    "key_id": "ck-7f3a",
    "nonce": "9c1f5e0a77d24b3e8a6f21d0c4b7e913",
    "seq": 6,
    "hash": "deadbeef...",
    "signature": "q3Zk0f..."}
<- {"status": "rejected", "session_id": "sess-001",
    "reason": "prefix_mismatch"}
\end{lstlisting}

\subsection{What is not in the wire format}

\begin{itemize}[leftmargin=2em]
\item \textbf{Transport layer} (HTTP, WebSocket, QUIC): out of scope.
The wire format is the JSON message content; the transport carries
it. A transport binding MAY recommend \texttt{Content-Type:
application/bridge-contract+json} for HTTP, but this is not
normative. The client authenticates the bridge via the transport
binding (e.g., TLS), which is out of scope here; the wire format
assumes an authenticated channel in the bridge-to-client direction.
\item \textbf{Streaming or batching}: each message is one request,
one response. A future revision MAY define a batch envelope.
\item \textbf{Bridge-to-client authentication}: the bridge is
trusted (A4); the client authenticates the bridge via the transport
binding (e.g., TLS), which is out of scope here.
\item \textbf{Operator actions} (\texttt{OPERATOR\_ADD\_KEY},
\texttt{OPERATOR\_TERMINATION}, \texttt{KEY\_REVOKED}):
not client-initiated. Operator actions travel on a separate operator
channel; the operator key is bound out-of-band (GAR \S 6.1), and the bridge
verifies the operator signature on each action before it commits (see the
operator signatures below). The reference front-ends give these actions
a request form in the same envelope (\texttt{revoke}, \texttt{add\_key},
\texttt{operator\_terminate}), so that the signature check is exercised
from a message; a deployment may carry them otherwise.
\item \textbf{Key exchange}: keys are registered at \texttt{attach}
and \texttt{add\_client}; there is no key-agreement message.
\end{itemize}

% ================================================================================
